% !TEX program = xelatex
\documentclass[12pt]{article}
\usepackage[paperwidth=612bp,paperheight=792bp,margin=0bp,noheadfoot]{geometry}
\usepackage{fontspec}
\usepackage{xeCJK}
\usepackage{tikz}
\pagestyle{empty}
\setlength{\parindent}{0pt}
\setlength{\parskip}{0pt}
\IfFontExistsTF{TeX Gyre Termes}{\setmainfont{TeX Gyre Termes}}{\setmainfont{Times New Roman}}
\IfFontExistsTF{Noto Serif CJK SC}{\setCJKmainfont{Noto Serif CJK SC}}{\setCJKmainfont{Songti SC}}
\newcommand{\QuestionTitle}{人机实验1 H-C6a：高祖本纪}
\newcommand{\DrawQuestionTitle}{%
\node[anchor=base west,inner sep=0pt,outer sep=0pt] at ([xshift=72bp,yshift=-34bp]current page.north west) {{\fontsize{14bp}{14bp}\selectfont\bfseries \QuestionTitle}};%
\node[anchor=base west,inner sep=0pt,outer sep=0pt] at ([xshift=72bp,yshift=-62bp]current page.north west) {{\fontsize{12bp}{12bp}\selectfont 用时：\underline{\hspace{80bp}}}};%
}
\begin{document}
% Auto-generated for 刘邦
\thispagestyle{empty}
\begin{tikzpicture}[remember picture,overlay]
\DrawQuestionTitle
\node[anchor=base west,inner sep=0pt,outer sep=0pt] at ([xshift=72bp,yshift=-84bp]current page.north west) {{\fontsize{12bp}{12bp}\selectfont 【8.9.1】 秦二世元年秋天，陈胜等在蕲县起义，到了陈县自立为王，号称“张楚”。}};
\node[anchor=base west,inner sep=0pt,outer sep=0pt] at ([xshift=72bp,yshift=-100bp]current page.north west) {{\fontsize{12bp}{12bp}\selectfont 【8.9.2】 各郡县都大多杀死长官，响应陈胜。}};
\node[anchor=base west,inner sep=0pt,outer sep=0pt] at ([xshift=72bp,yshift=-116bp]current page.north west) {{\fontsize{12bp}{12bp}\selectfont 【8.9.3】 沛县县令恐惧，想要以沛具响应陈胜，}};
\node[anchor=base west,inner sep=0pt,outer sep=0pt] at ([xshift=72bp,yshift=-132bp]current page.north west) {{\fontsize{12bp}{12bp}\selectfont 【8.9.4】 主吏萧何、狱掾曹参对他说：“您身为秦朝的官吏，如今要叛秦起事，率领沛}};
\node[anchor=base west,inner sep=0pt,outer sep=0pt] at ([xshift=72bp,yshift=-148bp]current page.north west) {{\fontsize{12bp}{12bp}\selectfont 县子弟，恐怕他们不愿听命。希望您召集逃亡在外面的人,可以得到几百人。利用这股力}};
\node[anchor=base west,inner sep=0pt,outer sep=0pt] at ([xshift=72bp,yshift=-164bp]current page.north west) {{\fontsize{12bp}{12bp}\selectfont 量胁持群众，群众不敢不听您的命令。”}};
\node[anchor=base west,inner sep=0pt,outer sep=0pt] at ([xshift=72bp,yshift=-180bp]current page.north west) {{\fontsize{12bp}{12bp}\selectfont 【8.9.5】 县令就派樊哙去召唤刘季，刘季的队伍已经近百人了。}};
\node[anchor=base west,inner sep=0pt,outer sep=0pt] at ([xshift=72bp,yshift=-196bp]current page.north west) {{\fontsize{12bp}{12bp}\selectfont 【8.10.1】 于是樊哙跟着刘季来到沛县。}};
\node[anchor=base west,inner sep=0pt,outer sep=0pt] at ([xshift=72bp,yshift=-212bp]current page.north west) {{\fontsize{12bp}{12bp}\selectfont 【8.10.2】 沛县县令又后悔了，恐怕刘季发生变故，就关闭城门，派人防守，（不让刘季}};
\node[anchor=base west,inner sep=0pt,outer sep=0pt] at ([xshift=72bp,yshift=-228bp]current page.north west) {{\fontsize{12bp}{12bp}\selectfont 进城，）打算杀掉萧何、曹参。}};
\node[anchor=base west,inner sep=0pt,outer sep=0pt] at ([xshift=72bp,yshift=-244bp]current page.north west) {{\fontsize{12bp}{12bp}\selectfont 【8.10.3】 萧何、曹参恐惧，翻过城墙依附刘季。}};
\node[anchor=base west,inner sep=0pt,outer sep=0pt] at ([xshift=72bp,yshift=-260bp]current page.north west) {{\fontsize{12bp}{12bp}\selectfont 【8.10.4】 刘季用帛写了一封信，}};
\node[anchor=base west,inner sep=0pt,outer sep=0pt] at ([xshift=72bp,yshift=-276bp]current page.north west) {{\fontsize{12bp}{12bp}\selectfont 【8.10.5】 射到城上，告诉沛县父老说：“天下苦于秦朝的暴政已经很久了。现在父老为}};
\node[anchor=base west,inner sep=0pt,outer sep=0pt] at ([xshift=72bp,yshift=-292bp]current page.north west) {{\fontsize{12bp}{12bp}\selectfont 沛令守城，但各国诸侯都已起事，（一旦城破，）就要屠戮沛县。如果沛县父老共同起来}};
\node[anchor=base west,inner sep=0pt,outer sep=0pt] at ([xshift=72bp,yshift=-308bp]current page.north west) {{\fontsize{12bp}{12bp}\selectfont 杀死沛令，选择子弟中可以立为首领的做领导，以响应诸侯军，那就能保全身家性命。不}};
\node[anchor=base west,inner sep=0pt,outer sep=0pt] at ([xshift=72bp,yshift=-324bp]current page.north west) {{\fontsize{12bp}{12bp}\selectfont 然的话，父子全遭杀害，死得毫无意义。”}};
\node[anchor=base west,inner sep=0pt,outer sep=0pt] at ([xshift=72bp,yshift=-340bp]current page.north west) {{\fontsize{12bp}{12bp}\selectfont 【8.10.6】 父老们就率领子弟共同杀了沛令，打开城门，迎接刘季，想让他做柿县县令。}};
\node[anchor=base west,inner sep=0pt,outer sep=0pt] at ([xshift=72bp,yshift=-356bp]current page.north west) {{\fontsize{12bp}{12bp}\selectfont 【8.10.7】 刘季说：“天下正在混乱当中，诸侯都已起事，如果推选的将领不胜任，就会}};
\node[anchor=base west,inner sep=0pt,outer sep=0pt] at ([xshift=72bp,yshift=-372bp]current page.north west) {{\fontsize{12bp}{12bp}\selectfont 一败涂地。我不是吝惜自己的生命，只怕才劣力薄，不能保全父兄子弟。这是件大事，希}};
\node[anchor=base west,inner sep=0pt,outer sep=0pt] at ([xshift=72bp,yshift=-388bp]current page.north west) {{\fontsize{12bp}{12bp}\selectfont 望另外共同推选一位能够胜任的人。”}};
\node[anchor=base west,inner sep=0pt,outer sep=0pt] at ([xshift=72bp,yshift=-404bp]current page.north west) {{\fontsize{12bp}{12bp}\selectfont 【8.10.8】 萧何、曹参等都是文官，看重身家性命，怕事情不成，秦朝会诛灭他们的全族}};
\node[anchor=base west,inner sep=0pt,outer sep=0pt] at ([xshift=72bp,yshift=-420bp]current page.north west) {{\fontsize{12bp}{12bp}\selectfont ，所以都推让刘季。}};
\node[anchor=base west,inner sep=0pt,outer sep=0pt] at ([xshift=72bp,yshift=-436bp]current page.north west) {{\fontsize{12bp}{12bp}\selectfont 【8.10.9】 父老们都说：“我们平时听到刘季许多奇异的事情，看来刘季是该显贵的。而}};
\node[anchor=base west,inner sep=0pt,outer sep=0pt] at ([xshift=72bp,yshift=-452bp]current page.north west) {{\fontsize{12bp}{12bp}\selectfont 且又经过占卜，没有比刘季更吉利的。”}};
\node[anchor=base west,inner sep=0pt,outer sep=0pt] at ([xshift=72bp,yshift=-468bp]current page.north west) {{\fontsize{12bp}{12bp}\selectfont 【8.10.10】 这时刘季再三谦让，大家都不敢担任，最后还是立刘季为沛公。}};
\node[anchor=base west,inner sep=0pt,outer sep=0pt] at ([xshift=72bp,yshift=-484bp]current page.north west) {{\fontsize{12bp}{12bp}\selectfont 【8.10.11】 在沛县衙门的庭院里祭祀黄帝和蚩尤，又用牲血衅鼓旗。}};
\node[anchor=base west,inner sep=0pt,outer sep=0pt] at ([xshift=72bp,yshift=-500bp]current page.north west) {{\fontsize{12bp}{12bp}\selectfont 【8.10.12】 旗子一律红色，因为刘季所杀蛇是白帝的儿子，杀蛇的是赤帝的儿子，所以}};
\node[anchor=base west,inner sep=0pt,outer sep=0pt] at ([xshift=72bp,yshift=-516bp]current page.north west) {{\fontsize{12bp}{12bp}\selectfont 崇尚赤色。}};
\node[anchor=base west,inner sep=0pt,outer sep=0pt] at ([xshift=72bp,yshift=-532bp]current page.north west) {{\fontsize{12bp}{12bp}\selectfont 【8.10.13】 于是少年子弟和有势的官吏，如萧何、曹参、樊哙等人，都为沛公征集兵员}};
\node[anchor=base west,inner sep=0pt,outer sep=0pt] at ([xshift=72bp,yshift=-548bp]current page.north west) {{\fontsize{12bp}{12bp}\selectfont ，集合了两三千人，攻打胡陵、方与，回军固守丰邑。}};
\node[anchor=base west,inner sep=0pt,outer sep=0pt] at ([xshift=72bp,yshift=-564bp]current page.north west) {{\fontsize{12bp}{12bp}\selectfont 【8.11.1】 秦二世二年，陈胜将领周章的军队西至戏水而还。}};
\node[anchor=base west,inner sep=0pt,outer sep=0pt] at ([xshift=72bp,yshift=-580bp]current page.north west) {{\fontsize{12bp}{12bp}\selectfont 【8.24.2】 汉高祖元年，秦地的百姓大失所望，然而心里恐惧，不敢不服从。}};
\node[anchor=base west,inner sep=0pt,outer sep=0pt] at ([xshift=72bp,yshift=-596bp]current page.north west) {{\fontsize{12bp}{12bp}\selectfont 【8.25.1】 项羽派人回去报告楚怀王。}};
\node[anchor=base west,inner sep=0pt,outer sep=0pt] at ([xshift=72bp,yshift=-612bp]current page.north west) {{\fontsize{12bp}{12bp}\selectfont 【8.25.2】 楚怀王说：“按照原来的约定办。”}};
\node[anchor=base west,inner sep=0pt,outer sep=0pt] at ([xshift=72bp,yshift=-628bp]current page.north west) {{\fontsize{12bp}{12bp}\selectfont 【8.25.3】 项羽怨恨楚怀王不肯让他与沛公一起西进入关，而派他北上救赵，}};
\node[anchor=base west,inner sep=0pt,outer sep=0pt] at ([xshift=72bp,yshift=-644bp]current page.north west) {{\fontsize{12bp}{12bp}\selectfont 【8.25.4】 在天下诸侯争夺称王关中的约定中落在后面，他就说：“怀王这个人，我家项}};
\node[anchor=base west,inner sep=0pt,outer sep=0pt] at ([xshift=72bp,yshift=-660bp]current page.north west) {{\fontsize{12bp}{12bp}\selectfont 梁所立，没有什么功劳，凭什么主持约定。本来安定天下的，是诸位将领和我项籍。”}};
\node[anchor=base west,inner sep=0pt,outer sep=0pt] at ([xshift=72bp,yshift=-676bp]current page.north west) {{\fontsize{12bp}{12bp}\selectfont 【8.25.5】 就假意推尊楚怀王为义帝，实际上不听从他的命令。}};
\node[anchor=base west,inner sep=0pt,outer sep=0pt] at ([xshift=72bp,yshift=-692bp]current page.north west) {{\fontsize{12bp}{12bp}\selectfont 【8.26.1】 正月，项羽自立为西楚霸王，在梁、楚地区的九个郡称王，建都彭城。}};
\node[anchor=base west,inner sep=0pt,outer sep=0pt] at ([xshift=72bp,yshift=-708bp]current page.north west) {{\fontsize{12bp}{12bp}\selectfont 【8.26.2】 背弃原来的约定，改立沛公为汉王，在巴、蜀、汉中称玉，建都南郑。}};
\node[anchor=base west,inner sep=0pt,outer sep=0pt] at ([xshift=72bp,yshift=-724bp]current page.north west) {{\fontsize{12bp}{12bp}\selectfont 【8.26.3】 把关中瓜分为三，封立秦朝的三个将领：章邯为雍王，建都废丘：司马欣为塞}};
\end{tikzpicture}
\null
\newpage
\thispagestyle{empty}
\begin{tikzpicture}[remember picture,overlay]
\DrawQuestionTitle
\node[anchor=base west,inner sep=0pt,outer sep=0pt] at ([xshift=72bp,yshift=-84bp]current page.north west) {{\fontsize{12bp}{12bp}\selectfont 王，建都栎阳；董翳为翟王，建都高奴。}};
\node[anchor=base west,inner sep=0pt,outer sep=0pt] at ([xshift=72bp,yshift=-100bp]current page.north west) {{\fontsize{12bp}{12bp}\selectfont 【8.26.4】 封楚将瑕丘申阳为河南王，建都洛阳。}};
\node[anchor=base west,inner sep=0pt,outer sep=0pt] at ([xshift=72bp,yshift=-116bp]current page.north west) {{\fontsize{12bp}{12bp}\selectfont 【8.26.5】 封赵将司马印为殷王，建都朝歌。}};
\node[anchor=base west,inner sep=0pt,outer sep=0pt] at ([xshift=72bp,yshift=-132bp]current page.north west) {{\fontsize{12bp}{12bp}\selectfont 【8.26.6】 赵王歇迁徙代地称王。}};
\node[anchor=base west,inner sep=0pt,outer sep=0pt] at ([xshift=72bp,yshift=-148bp]current page.north west) {{\fontsize{12bp}{12bp}\selectfont 【8.26.7】 封赵将张耳为常山王，建都襄国。}};
\node[anchor=base west,inner sep=0pt,outer sep=0pt] at ([xshift=72bp,yshift=-164bp]current page.north west) {{\fontsize{12bp}{12bp}\selectfont 【8.26.8】 封当阳君黥布为九江王，建都六县。}};
\node[anchor=base west,inner sep=0pt,outer sep=0pt] at ([xshift=72bp,yshift=-180bp]current page.north west) {{\fontsize{12bp}{12bp}\selectfont 【8.26.9】 封楚怀王柱国共敖为临江王，建都江陵。}};
\node[anchor=base west,inner sep=0pt,outer sep=0pt] at ([xshift=72bp,yshift=-196bp]current page.north west) {{\fontsize{12bp}{12bp}\selectfont 【8.26.10】 封番郡吴芮为衡山王，建都邾县。}};
\node[anchor=base west,inner sep=0pt,outer sep=0pt] at ([xshift=72bp,yshift=-212bp]current page.north west) {{\fontsize{12bp}{12bp}\selectfont 【8.26.11】 封燕将臧荼为燕王，建都蓟县。}};
\node[anchor=base west,inner sep=0pt,outer sep=0pt] at ([xshift=72bp,yshift=-228bp]current page.north west) {{\fontsize{12bp}{12bp}\selectfont 【8.26.12】 原来的燕王韩广迁徙辽东称王。}};
\node[anchor=base west,inner sep=0pt,outer sep=0pt] at ([xshift=72bp,yshift=-244bp]current page.north west) {{\fontsize{12bp}{12bp}\selectfont 【8.26.13】 韩广不服从，臧荼攻杀韩广于无终。}};
\node[anchor=base west,inner sep=0pt,outer sep=0pt] at ([xshift=72bp,yshift=-260bp]current page.north west) {{\fontsize{12bp}{12bp}\selectfont 【8.26.14】 封成安君陈余河间三县，住在南皮。}};
\node[anchor=base west,inner sep=0pt,outer sep=0pt] at ([xshift=72bp,yshift=-276bp]current page.north west) {{\fontsize{12bp}{12bp}\selectfont 【8.26.15】 封给梅鋗十万户。}};
\node[anchor=base west,inner sep=0pt,outer sep=0pt] at ([xshift=72bp,yshift=-292bp]current page.north west) {{\fontsize{12bp}{12bp}\selectfont 【8.27.1】 四月，在项羽旌麾之下罢兵散归，诸侯各自回到封国。}};
\node[anchor=base west,inner sep=0pt,outer sep=0pt] at ([xshift=72bp,yshift=-308bp]current page.north west) {{\fontsize{12bp}{12bp}\selectfont 【8.29.1】 八月，汉王用韩信的计策，从故道回军，袭击雍王章邯。}};
\node[anchor=base west,inner sep=0pt,outer sep=0pt] at ([xshift=72bp,yshift=-324bp]current page.north west) {{\fontsize{12bp}{12bp}\selectfont 【8.29.2】 章邯在陈仓迎击汉军，雍王兵败退走，在好畤停下来接战，又失败了，逃到废}};
\node[anchor=base west,inner sep=0pt,outer sep=0pt] at ([xshift=72bp,yshift=-340bp]current page.north west) {{\fontsize{12bp}{12bp}\selectfont 丘。}};
\node[anchor=base west,inner sep=0pt,outer sep=0pt] at ([xshift=72bp,yshift=-356bp]current page.north west) {{\fontsize{12bp}{12bp}\selectfont 【8.29.3】 汉王随即平定了雍地。}};
\node[anchor=base west,inner sep=0pt,outer sep=0pt] at ([xshift=72bp,yshift=-372bp]current page.north west) {{\fontsize{12bp}{12bp}\selectfont 【8.29.4】 向东到达咸阳，率军围困雍王于废丘，而派遣将领攻占了陇西、北地、上郡。}};
\node[anchor=base west,inner sep=0pt,outer sep=0pt] at ([xshift=72bp,yshift=-388bp]current page.north west) {{\fontsize{12bp}{12bp}\selectfont 【8.29.5】 派将军薛欧、王吸出武关，借助王陵驻扎在南阳的兵力，迎接太公、吕后于沛}};
\node[anchor=base west,inner sep=0pt,outer sep=0pt] at ([xshift=72bp,yshift=-404bp]current page.north west) {{\fontsize{12bp}{12bp}\selectfont 县。}};
\node[anchor=base west,inner sep=0pt,outer sep=0pt] at ([xshift=72bp,yshift=-420bp]current page.north west) {{\fontsize{12bp}{12bp}\selectfont 【8.29.6】 楚听到这一消息，出兵在阳夏阻挡，汉军不能前进。}};
\node[anchor=base west,inner sep=0pt,outer sep=0pt] at ([xshift=72bp,yshift=-436bp]current page.north west) {{\fontsize{12bp}{12bp}\selectfont 【8.29.7】 楚让原吴县县令郑昌为韩王，抵抗汉军。}};
\node[anchor=base west,inner sep=0pt,outer sep=0pt] at ([xshift=72bp,yshift=-452bp]current page.north west) {{\fontsize{12bp}{12bp}\selectfont 【8.30.1】 二年，汉王东出略取城邑，塞王司马欣、翟王董翳、河南王申阳都投降了。}};
\end{tikzpicture}
\null
\end{document}

